\section{Cubic extensions on the blow-up}\label{blowup}

Let $I\subseteq\m$ be as in Theorem C, or let $I=\m$ when $\operatorname{gr}_{\m}R$ is a hypersurface in four variables. In the latter case $R$ is automatically Cohen--Macaulay of dimension three. Indeed a homogeneous system of parameters in the graded hypersurface is a regular sequence. Its lifts are regular in $R$ by induction, using $\operatorname{gr}(R/aR)=\operatorname{gr}(R)/(a^*)$ whenever the initial form $a^*$ is a nonzerodivisor. Write
\[
 X=\operatorname{Bl}_I(\Spec R),\quad D=V(I\mathcal O_X),\quad
 L=I\mathcal O_X=\mathcal O_X(-D),\quad T=X\setminus D.
\]
The exceptional divisor is effective Cartier, $L$ is relatively ample, and $L|_D=\mathcal O_D(1)$ \cite{Stacks}. We identify $T$ with $\Spec R\setminus V(I)$. Since $I$ contains a regular sequence of length three, the local cohomology sequence for $V(I)$ gives $\Gamma(T,\mathcal O_T)=R$. We write $L^n\mathcal E$ for $L^{\otimes n}\otimes\mathcal E$, allowing negative $n$.

\subsection{Multiplication and the obstruction to splitting}
Let $U$ be the matrix presenting $S/R$ and write $s=b+2$. Dividing each column by a local equation for $D$ gives a bundle sequence
\begin{equation}\label{bundlec}
 0\longrightarrow (L^{-1})^b\longrightarrow\mathcal O_X^s
       \xrightarrow{\rho} V\longrightarrow0 .
\end{equation}
Indeed, the divided matrix on the closed fiber is $\overline U_I$. The rank-drop locus is closed in $X$ and proper over the local base. Its image misses the closed point, so it is empty. Thus \eqref{bundlec} is locally split. The bundle $V$ has rank two, $\det V=L^b$, and $V|_T=(S/R)|_T$. The kernel of $U:R^b\to R^s$ is supported in $V(I)$ and lies in a free module. A regular element of $I$ therefore forces the kernel to vanish. In particular $S/R$ has projective dimension one and $\pd_R S\le1$.

The algebra $S|_T$ is locally free of rank three. Formula~\eqref{perfect-trace} settles the case $3\in R^\times$. Henceforth the residue characteristic is three, so $2$ is a unit. Define
\[
 M=\mathcal Hom(\Sym^2 V,V),\qquad Q=\Sym^3(V^\vee)\otimes\det V.
\]
There is an exact sequence
\begin{equation}\label{cubicsequence}
 0\longrightarrow V^\vee\xrightarrow{i}M\xrightarrow{j}Q\longrightarrow0,
 \qquad j(\mu)(v)=v\wedge\mu(v,v),
\end{equation}
where $i(\ell)(v,w)=\ell(v)w+\ell(w)v$. In a basis write
$\mu_{11}=(a,b'),\ \mu_{12}=(c,d),\ \mu_{22}=(h,f)$.
The four coefficients of $j(\mu)$ are $b',2d-a,f-2c,-h$. They give surjectivity and identify the kernel with $i(V^\vee)$.

A local linear section $\sigma$ of $S|_T\to V|_T$ gives
$\mu_\sigma(v,w)=\overline{\sigma(v)\sigma(w)}$. Changing $\sigma$ by a linear form $\ell$ changes $\mu_\sigma$ by $i(\ell)$. Hence $j(\mu_\sigma)$ is a global binary cubic on $T$ \cite[Theorem~2.1]{Wood}.

\begin{lemma}\label{cubic-lifting}
This cubic extends to $\Phi\in H^0(X,Q)$. If $\Phi$ lifts to $H^0(X,M)$, then $R\subseteq S$ splits.
\end{lemma}
\begin{proof}
Lift the $s$ generators of $S/R$ to $S$ to obtain a map $h:\mathcal O_X^s\to S\otimes_R\mathcal O_X$. The quotient $S\otimes_R\mathcal O_X\to V$ is well defined since the image of $U$ is contained in that of the divided matrix. Denote it by $g$. The polynomial
\[
 z\longmapsto\rho(z)\wedge g(h(z)^2)
\]
on $\mathcal O_X^s$ restricts on $T$ to the pullback of the binary cubic. Locally split \eqref{bundlec}. All coefficients containing a kernel variable vanish on $T$, as polynomial coefficients over the local coordinate ring. A local equation of $D$ is a nonzerodivisor, so these coefficients vanish on $X$ as well. The remaining coefficients give the desired section, uniquely.

If $j(\mu)=\Phi$, write $\mu_{\sigma_i}-\mu=i(\ell_i)$ on opens of $T$ with local sections $\sigma_i$. On overlaps, injectivity of $i$ gives $\sigma_i-\sigma_j=\ell_i-\ell_j$. The sections $\sigma_i-\ell_i$ glue and split the unit sequence on $T$. Evaluate the resulting projection $S|_T\to\mathcal O_T$ on elements of $S$ and use $\Gamma(T,\mathcal O_T)=R$.
\end{proof}

Rank-two duality gives $L^{2b}Q=\Sym^3 V$ and
$L^{2b}M\simeq V\otimes\Sym^2 V$. The latter identification uses $2\in R^\times$ to polarize quadratic forms. Under these identifications $L^{2b}j$ is symmetric multiplication up to a sign. Thus a cubic monomial in the global generators of $V$ has a global lift through this twisted map.

\subsection{A vanishing in residue characteristic three}
\begin{lemma}\label{frobcube}
Let $Y$ be a scheme over a field of characteristic three with an invertible sheaf $\mathcal O_Y(1)$. Suppose
$H^0(Y,\mathcal O_Y(t))=0$ for $t<0$ and
$H^1(Y,\mathcal O_Y(t))=0$ for every $t$. If
\[
 0\longrightarrow\mathcal O_Y(-1)^b\longrightarrow
 \mathcal O_Y^{b+2}\longrightarrow E\longrightarrow0
\]
is a bundle sequence, then $H^0(Y,(\Sym^3 E)(t))=0$ for $t<-b$.
\end{lemma}
\begin{proof}
For a rank-two bundle in characteristic three the absolute Frobenius $F_Y$ gives an exact sequence
\begin{equation}\label{frob-cube-sequence}
 0\longrightarrow F_Y^*E\longrightarrow\Sym^3 E
 \longrightarrow E\otimes\det E\longrightarrow0 .
\end{equation}
In a basis $x,y$ the first arrow has image spanned by $x^3,y^3$. To see that the quotient is canonical, take the degree-three derivative
$\Sym^3 E\to E\otimes\Sym^2 E$, defined by the Leibniz rule. Its image is the image of the injection
\[
 E\otimes\det E\longrightarrow E\otimes\Sym^2 E,\qquad
 a\otimes(u\wedge v)\longmapsto u\otimes(av)-v\otimes(au).
\]
Indeed, the derivatives of $x^2y,xy^2$ give its two basis vectors up to sign, and the derivatives of the cubes vanish. This coefficient calculation works on nonreduced schemes.

The given resolution gives $H^0(Y,E(j))=0$ for $j<0$. Its Frobenius pullback is exact by local splitting and replaces the left-hand twist by $-3$. Thus $H^0(Y,(F_Y^*E)(t))=0$ for $t<0$ as well. Since $\det E=\mathcal O_Y(b)$, the two outer terms of \eqref{frob-cube-sequence} twisted by $t$ have zero global sections for $t<-b$.
\end{proof}

\begin{proof}[Proof of Theorem C]
In case (i) set $B=R/I$. Since $I$ is generated by a regular sequence, the conormal module is free of rank three and
\[
 \operatorname{gr}_I R=\Sym_B(I/I^2)\simeq B[X_0,X_1,X_2].
\]
Thus $D=\PP^2_B$ and its closed fiber is $Y=\PP^2_k$. In case (ii), $D=Y=V(f_2)\subset\PP^3_k$. Both closed fibers satisfy the line-bundle hypotheses of Lemma~\ref{frobcube}. For the quadric this follows by twisting $0\to\mathcal O_{\PP^3}(-2)\xrightarrow{f_2}\mathcal O_{\PP^3}\to\mathcal O_D\to0$. No smoothness hypothesis enters.

Let $E=V|_Y$. Restriction of \eqref{bundlec} gives precisely the presentation in Lemma~\ref{frobcube}, and
\[
 (L^n Q)|_Y\simeq (\Sym^3 E)(n-2b).
\]
For $0\le n<b$, the lemma gives zero global sections on $Y$. In case (i) the same vanishing holds on all of $D$. Here is the argument when $B$ is arbitrary. Put $\mathfrak a=\m/I$ and $\mathcal H=(L^nQ)|_D$. This bundle is flat over $B$, so its $j$th graded piece for the $\mathfrak a$-adic filtration is
$(\mathfrak a^j/\mathfrak a^{j+1})\otimes_k\mathcal H|_Y$. None of these pieces has a nonzero global section. A global section of $\mathcal H$ therefore belongs to $\mathfrak a^j\mathcal H$ for every $j$. Krull intersection makes its germ zero at each point over the closed point of $B$. Its support is proper over $B$, hence empty. This also explains why characteristic three is needed only on the closed fiber.

Use the restriction sequences for the Cartier divisor $D$ successively. The vanishing just proved places $\Phi$ in $H^0(X,L^bQ)$. To lift it, note that in both cases
\[
 H^1(D,\mathcal O_D(j))=0\quad(j\ge0),\qquad
 H^2(D,\mathcal O_D(j))=0\quad(j\ge-1).
\]
For $\PP^2_B$ these are the cohomology formulas for projective space over a ring \cite{Stacks}. For the quadric they follow from the same hypersurface sequence as above. Since the base is affine, relative Serre vanishing followed by descending induction in
\[
 0\longrightarrow L^{j+1}\longrightarrow L^j
       \longrightarrow\mathcal O_D(j)\longrightarrow0
\]
gives $H^1(X,\mathcal O_X)=H^2(X,L^{-1})=0$. Sequence \eqref{bundlec} now gives $H^1(X,V)=0$. Finally $L^b V^\vee\simeq V$, so \eqref{cubicsequence} twisted by $L^b$ lifts $\Phi$ to $H^0(X,L^bM)$. Lemma~\ref{cubic-lifting} gives the retraction. The closed-fiber matrix depends only on $U$ modulo $\m I$.
\end{proof}

\begin{corollary}\label{toeplitz}
Let $I=(a_0,a_1,a_2)$ be generated by a regular sequence in a local ring. For each $b\ge1$, take the $(b+2)\times b$ matrix whose $j$th column has $a_0,a_1,a_2$ in successive positions $j,j+1,j+2$ and zeros elsewhere. Every module-finite extension with unit quotient presented by this matrix, or by any perturbation of it in $\m I$, splits.
\end{corollary}
\begin{proof}
At a geometric point $[x_0:x_1:x_2]$ of $\PP^2_k$, the matrix sends a polynomial of degree $<b$ to its product with $x_0+x_1 Z+x_2 Z^2$. This polynomial is nonzero, so multiplication is injective. Apply Theorem C(i).
\end{proof}

For instance take $R=\mathbf Z_3[[x,y,z]]/(3^e+x^e+y^e+z^e)$ for any $e\ge2$. The ideal $(3,x,y)$ is a parameter ideal since its quotient is $\mathbf F_3[[z]]/(z^e)$, and $R$ is Cohen--Macaulay. Thus Corollary~\ref{toeplitz} applies for every $b$, including $b\ge4$, at arbitrarily large multiplicity. More generally, for any fixed triple as above, full column rank on $\PP^2_k$ defines an open subset of the affine space of coefficient matrices because the rank-drop incidence is closed and its projection is proper. The displayed matrices give a $k$-point of this open over every residue field.
